%!TEX root = Love602.tex
606 CONICAL SHELL CH. XXIV A %\section{Article 352}

in the case of lateral forces. The equations of the system fall into two
sets. One set arises from the equations of equilibrium, which become

\[\begin{aligned}
& \frac{d}{d r}\left(t_1 r\right) \sin \gamma-n s_2-t_2 \sin \gamma=0, \\
& \frac{d}{d r}\left(s_1 r\right) \sin \gamma-n t_2-s_2 \sin \gamma-n_2 \cos \gamma=0, \\
& \frac{d}{d r}\left(n_1 r\right) \sin \gamma+n n_2+t_2 \cos \gamma=0, \\
& \frac{d}{d r}\left(h_1 r\right) \sin \gamma+n g_2+h_1 \sin \gamma+n_2 r \sin \gamma=0, \\
& \frac{\omega}{d r}\left(g_1 r\right) \sin \gamma-n h_1-g_2 \sin \gamma-n_1 r \sin \gamma=0, \\
& -h_1 \cos \gamma+\left(s_1+s_2\right) r \sin \gamma=0 .
\end{aligned}\]

The other set arises from the equations expressing the stress-resultants
and stress-couples in terms of the displacement. The equations of this set are
\[\begin{aligned}
& t_1=\frac{d U}{d r}+\sigma\left(\frac{U}{r}+n V-W \cos \gamma\right. \\
& r \sin \gamma \\
& t_2=\sigma \frac{d U}{d r}+\left(\frac{U}{r}+\frac{n V-W \cos \gamma}{r \sin \gamma}\right), \\
& s_1-s_2=(1-\sigma)\left(\frac{d V}{d r}-\frac{V}{r}-\frac{n U}{r \sin \gamma}\right), \\
& g_1=-\frac{h^2}{3}\left\{\frac{d^2 W}{d r^2}+\sigma\left(\frac{n V \cos \gamma-n^2 W}{r^2 \sin ^2 \gamma}+\frac{1}{r} \frac{d W}{d r}\right)\right\} \\
& g_2=-\frac{1}{2}\left\{\sigma \frac{d^2 W}{d r^2}+\left(\frac{n V \cos \gamma-n^2 W}{r^2 \sin ^2 \gamma}+\frac{1}{r} \frac{d W}{d r}\right)\right\} \\
& h_1=\frac{h^2}{3}(1-\sigma) \frac{d}{d r}\left(\frac{V \cos \gamma-n W}{r \sin \gamma}\right) .
\end{aligned}\]

It will be observed that the 6th of equations (145) and the 3rd of
equations (146) are equivalent to the formulæ for \(S_1, S_2\) in
Article 330.

We shall simplify the equations (146) a little by introducing instead of
\(U, V, W\) the quantities \(\xi, \eta, \zeta\) defined by the equations

\[\xi=\frac{U}{r}, \quad=\frac{V \cos \gamma-n W}{r \sin \gamma}, \zeta=\frac{W}{r} .\]

Then the solution is effected in three stages. The first stage, marked
(a), consists in eliminating \(n_1, n_2, s_2\) by means of the 4th, 5th,
and 6th of equations (145), so as to obtain three equations containing
\(t_1, t_2, s_1, g_1, g_2, h_1\), and transforming these, by means of
(146) and (147), into three equations for determining
\(\xi, \eta, \zeta\). The second stage, marked (b), consists in assuming
for \(\xi, \eta, \zeta\) such forms as \(e^{m a} X\), where

\[\frac{x^2}{2}=\frac{r}{h},\]

%Article 352. EDGE-EFFECT. UNSYMMETRICAL CONDITIONS 607

and obtaining the equations which the quantities of the type \(X\) must
satisfy. The third stage, marked (c), consists in assuming the
quantities of the type \(X\) to be expressible in the form

\[z^a\left(c_0+c_1 z+c_2 z^2+\ldots\right),\]
where \(z=x^{-1}\), and determining the coefficients
\(m, a, c_0, c_1, \ldots\), in so far as these are not arbitrary.

\subsection{Formation of the equations.(a)}

In the first three of equations (145) we substitute for
\(s_2, n_1, n_2\) from the last three, obtaining the equations

\[\begin{aligned}
& \frac{d}{d r}\left(t_1 r\right)+\frac{n}{\sin \gamma}\left(s_1-\frac{h_1}{r \tan \gamma}\right)-t_2=0, \\
& \left\{\frac{d}{d r}\left(s_1 r\right)+s_1\right\} \sin \gamma+\frac{\cos \gamma}{r} \frac{d}{d r}\left(h_1 r\right)-n\left(t_2-\frac{g_2}{r \tan \gamma}\right)=0, \\
& \frac{d^2}{d r^2}\left(g_1 r\right) \sin \gamma-\frac{2 n}{r} \frac{d}{d r}\left(h_1 r\right)-\frac{d g_2}{d r} \sin \gamma-\frac{n^2 g_2}{r \sin \gamma}+t_2 \cos \gamma=0 .
\end{aligned}\]

We eliminate \(U, V, W\) from equations (146) by using equations (147),
obtaining the equations

\[\begin{aligned}
& t_1=\frac{d}{d r}(r \xi)+\sigma\left(\xi+n \eta \sec \gamma+\frac{n^2-\cos ^2 \gamma}{\sin \gamma \cos \gamma} \zeta\right), \\
& t_2=\sigma \frac{d}{d r}(r \xi)+\left(\xi+n \eta \sec \gamma+\frac{n^2-\cos ^2}{\sin \gamma \cos \gamma} \zeta\right), \\
& \quad \frac{1-\sigma}{2}\left\{-\frac{n}{\sin \gamma} \xi+\left(r \tan \gamma+\frac{h^2}{3 r \tan \gamma}\right) \frac{d \eta}{d r}+n r \frac{d \zeta}{d r} \sec \gamma\right\}, \\
& g_1: \quad \frac{h^2}{3}\left[\frac{d^2}{d r^2}(r \zeta)+\sigma\left\{\begin{array}{c}
n \eta \\
r \sin \gamma
\end{array}+\frac{1}{r} \frac{d}{d r}(r \zeta)\right\}\right], \\
& g_2=-\frac{h^2}{3}\left[\sigma \frac{d^2}{d r^2}\left(r()+\left\{\frac{n \eta}{r \sin \gamma}+\frac{1}{r} \frac{d}{d r}(r \zeta)\right\}\right],\right. \\
& h_1=\frac{h^2}{3}(1-\sigma) \frac{d \eta}{d r} .
\end{aligned}\]

len we put

\[r=\rho h,\]

Id substitute from (149) in (148) obtaining three equations from which
\(h\) is explicitly sent. The condition that \(h\) is small is now
replaced by the condition that a solution is quired for large values of
\(p\). A little simplification is effected by putting

\[\xi_1=\xi \rho, \quad \eta_1=\eta \rho, \quad \zeta_1=\zeta \rho .\]

Then the three equations become

\[
\left(\frac{d^2 \xi_1}{d \rho^2}+\frac{d \xi_1}{d \rho}-\frac{\xi_1}{\rho}\right) \sin \gamma-\frac{n^2(1-\sigma)}{2 \sin \gamma} \frac{\xi_1}{\rho}+
\]
\[
\sigma n \tan \gamma \frac{d \eta_1}{d \rho}+n \frac{1-\sigma}{2} \tan \gamma\left(\frac{d \eta_1}{d \rho}-\frac{\eta_1}{\rho}\right)
\]
\[
\frac{1-\sigma}{6 \tan \gamma} \frac{1}{\rho^2}\left(\frac{d \eta_1}{d \rho}-\frac{\eta_1}{\rho}\right)-n \tan \gamma \frac{\eta_1}{\rho}+\left(n^2 \frac{1+\sigma}{2 \cos \gamma}-\sigma \cos \gamma\right) \frac{d \zeta_1}{d \rho}
\]
\[
-n^2\left(\frac{3-\sigma}{2 \cos \gamma}-\cos \gamma\right) \frac{\zeta_1}{\rho}=0,
\]

608 CONICAL SHELL CH. XXIV A

(153)
\[\begin{array}{r}
-n \frac{1+\sigma}{2} \frac{d \xi_1}{d \rho}-n \frac{3-\sigma}{2} \frac{\xi_1}{\rho}+
\frac{1-\sigma}{2} \frac{\sin ^2 \gamma}{\cos \gamma}\left(\rho \frac{d^2 \eta_1}{d \rho^2}+\frac{d \eta_1}{d \rho}-\frac{\eta_1}{\rho}\right)+ \\
\frac{1-\sigma}{2} \cos \gamma\left(\frac{1}{\rho} \frac{d^2 \eta_1}{d \rho^2}-\frac{1}{\rho^2} \frac{d \eta_1}{d \rho}+\frac{\eta_1}{\rho^3}\right) \\
-\frac{n^2}{\cos \gamma} \frac{\eta_1}{\rho}-\frac{n^2 \cos \gamma}{3 \sin ^2 \gamma} \frac{\eta_1}{\rho^3}+\frac{1-\sigma}{2} n \tan \gamma\left(\rho \frac{d^2 \zeta_1}{d \rho^2}+\frac{d \zeta_1}{d \rho}-\frac{\zeta_1}{\rho}\right)+ \\
-\frac{\sigma n}{3 \tan \gamma} \frac{1}{\rho} \frac{d^2 \zeta_1}{d \rho^2}-\frac{n}{3 \tan \gamma} \frac{1}{\rho^2} \frac{d \zeta_1}{d \rho} \\
-n \frac{n^2-\cos 2 \gamma}{\sin \gamma \cos \gamma} \frac{\zeta_1}{\rho}=0, \ldots
\end{array}
\]
and
\[\begin{aligned}
\sigma \cos \gamma \frac{d \xi_1}{d \rho}+\frac{\xi_1}{\rho} \cos \gamma-n \frac{2-\sigma}{3} \frac{1}{\rho} \frac{d^2 \eta_1}{d \rho^2}+\frac{n}{\rho^2} \frac{d \eta_1}{d \rho}+n \frac{\eta_1}{\rho}+\left(\frac{n^3}{3 \sin ^2 \gamma}-\frac{4 n}{3}\right) \frac{\eta_1}{\rho^3} \\
-\frac{\sin \gamma}{3}\left(\rho \frac{d^4 \zeta_1}{d \rho^4}+2 \frac{d^3 \zeta_1}{d \rho^3}-\frac{1}{\rho} \frac{d^2 \zeta_1}{d \rho^2}+\frac{1}{\rho^2} \frac{d \zeta_1}{d \rho}\right)+\frac{\sigma n^2}{3 \sin \gamma} \frac{1}{\rho} \frac{d^2 \zeta_1}{d \rho^2}+\frac{n^2}{3 \sin \gamma} \frac{1}{\rho^2} \frac{d \zeta_1}{d \rho}+ \\
\frac{n^2-\cos ^2 \gamma}{\sin \gamma} \frac{\zeta_1}{\rho}=0 .
\end{aligned}\]

\subsection{Preparation for solution (b)}

In the three equations (152), (153), (154) we change the independent
variable by putting

\[\rho=\frac{1}{2} x^2 .\]

It will prove convenient presently to have \(x \xi_1\) as a dependent
variable instead of \(\xi_1\), so we put
\[\xi_2=x \xi_1,\]
and write the three equations in the forms
\[
\frac{\sin \gamma}{2}\left(\frac{d^2 \xi_2}{d x^2}-\frac{1}{x} \frac{d \xi_5}{d x}-\frac{3}{x^2} \xi_2\right)-n^2 \frac{1-\sigma}{\sin \gamma} \frac{\xi_2}{x^2}+
\]
\[
+n \frac{1+\sigma}{2} \frac{d \eta_1}{d x} \tan \gamma-n(3-\sigma) \frac{\eta_1}{x} \tan \gamma +
\]
\[
 -2 n \frac{1-\sigma}{3 \tan \gamma}\left(\frac{1}{x^4} \frac{d \eta_1}{d x}-\frac{2}{x^6} \eta_1\right)+\left(n^2 \frac{1+\sigma}{2 \cos \gamma}-\sigma \cos \gamma\right) \frac{d \zeta_1}{d x}+
\]
\[
-\left(n^2 \frac{3-\sigma}{\cos \gamma}-2 \cos \gamma\right) \frac{\zeta_1}{x}=0,
\]

\[\begin{aligned}
& -n \frac{1+\sigma}{2} \frac{1}{x^2} \frac{d \xi_2}{d x}-n \frac{5-3 \sigma}{2} \frac{\xi_2}{x^3}+\frac{1-\sigma}{4} \frac{\sin ^2 \gamma}{\cos \gamma}\left(\frac{d^2 \eta_1}{d x^2}+\frac{1}{x} \frac{d \eta_1}{d x}-\frac{4}{x^2} \eta_1\right) \\
& +(1-\sigma) \cos \gamma\left(\frac{1}{x^4} \frac{d^2 \eta_1}{d x^2}-\frac{3}{x^6} \frac{d \eta_1}{d x}+\frac{4}{x^6} \eta_1\right)-\frac{2 n^2}{\cos \gamma} \frac{\eta_1}{x^3}-\frac{8 n^3 \cos \gamma}{3 \sin ^2 \gamma} \frac{\eta_1}{x^6} \\
& +\frac{1-\sigma}{4} n \tan \gamma\left(\frac{d^2 \zeta_1}{d x^2}+\frac{1}{x} \frac{d \zeta_1}{d x}-\frac{4}{x^2} \zeta_1\right)-\frac{2 \sigma n}{3 \tan \gamma}\left(\frac{1}{x^4} \frac{d^2 \zeta_1}{d x^2}-\frac{1}{x^5} \frac{d \zeta_1}{d x}\right) \\
& -\frac{4 n}{3 \tan \gamma} \frac{1}{x^5} \frac{d \zeta_1}{d x}-2 n \frac{n^2-\cos 2}{d x} \frac{\zeta_1}{\sin \gamma \cos \gamma}=0,
\end{aligned}\]
and

\[\begin{array}{c}
\sigma \frac{d \xi_2}{d x} \cos \gamma+\frac{\psi-\sigma}{x} \xi_2 \cos \gamma-2 n \frac{2-\sigma}{3} \frac{1}{x^2} \frac{d^2 \eta_1}{d x^2}+2 n \frac{8-\sigma}{3} \frac{1}{x^3} \frac{d \eta_1}{d x}+\\
\left(\frac{8 n^3}{3 \sin ^2 \gamma}-\frac{32}{3} n\right) \frac{\eta_1}{x^4}+2 n \eta_1 \\
-\frac{\sin \gamma}{6}\left(\frac{d^4 \zeta_1}{d x^4}-\frac{2}{x} \frac{d^3 \zeta_1}{d x^3}-\frac{1}{x^2} \frac{d^2 \zeta_1}{d x^2}+\frac{9}{x^3} \frac{d \zeta_1}{d x}\right)+\frac{2 \sigma n^2}{3 \sin ^2 \gamma}\left(\frac{1}{x^2} \frac{d^2 \zeta_1}{d x^3}-\frac{1}{x^3} \frac{d \zeta_1}{d x}\right) \\
+\frac{4 n^2}{3 \sin \gamma} \frac{1}{x^3} \frac{d \zeta_1}{d x}+2 \frac{n^2-\cos ^2 \gamma}{\sin \gamma} \zeta_1=0 . \quad \ldots . .(159)
\end{array}\]

%352. EDGE-EFFECT. UNSYMMETRICAL CONDITIONS 609

Equations (157), (158), (159) are three linear differential equations
for the determination of \(\xi_2, \eta_1, \zeta_1\) as functions of
\(x\), and we seek solutions of a certain type answering to the normal
integrals of a single linear differential equation. The first step is to
put

\[\xi_2=X e^{m x}, \quad \eta_1=Y e^{m x}, \quad \zeta_1=Z e^{m x},\]
where \(m\) is at present undetermined. Then \(X, Y, Z\) satisfy the
three equations
\[
\begin{array}{c}
\frac{\sin \gamma}{2}\left\{\left(\frac{d^2 X}{d x^2}+2 m \frac{d X}{d x}+m^2 X\right)-\frac{1}{x}\left(\frac{d X}{d x}+m X\right)-\frac{3}{x^2} X\right\}
-n^2 \frac{1-\sigma}{\sin \gamma} \frac{X}{x^2} \\
+n \frac{1+\sigma}{2} \tan \gamma\left(\frac{d Y}{d x}+m Y\right)\\
-n(3-\sigma) \tan \gamma \frac{Y}{x}-2 n \frac{1-\sigma}{3 \tan \gamma}\left\{\frac{1}{x^4}\left(\frac{d Y}{d x}+m Y\right)-\frac{2}{x^6} Y\right\} \\
+\left(n^2 \frac{1+\sigma}{2 \cos \gamma}-\sigma \cos \gamma\right)\left(\frac{d Z}{d x}+m Z\right)-\left(n^2 \frac{3-\sigma}{\cos \gamma}-2 \cos \gamma\right) \frac{Z}{x}=0, \ldots \ldots(161) \\
-n \frac{1+\sigma}{2} \frac{1}{x^2}\left(\frac{d X}{d x}+m X\right)-n \frac{5-3 \sigma}{2} \frac{X}{x^3} \\
+\frac{1-\sigma}{4} \frac{\sin ^2 \gamma}{\cos \gamma}\left\{\left(\frac{d^2 Y}{d x^2}+2 m \frac{d Y}{d x}+m^2 Y\right)+\frac{1}{x}\left(\frac{d Y}{d x}+m Y\right)-\frac{4}{x^3} Y\right\} \\
+(1-\sigma) \cos \gamma\left\{\frac{1}{x^4}\left(\frac{d^2 Y}{d x^2}+2 m \frac{d Y}{d x}+m^2 Y\right)-\frac{3}{x^6}\left(\frac{d Y}{d x}+m Y\right)+\frac{4}{x^3} Y\right\} \\
-\frac{2 n^2}{\cos \gamma} \frac{Y}{x^2}-\frac{8 n^3 \cos \gamma}{3 \sin ^2 \gamma} \frac{Y}{x^3} \\
+\frac{1-\sigma}{4} n \tan \gamma\left\{\left(\frac{d^2 Z}{d x^2}+2 m \frac{d Z}{d x}+m^2 Z\right)+\frac{1}{x}\left(\frac{d Z}{d x}+m Z\right)-\frac{4}{x^2} Z\right\} \\
-\frac{2 \sigma n}{3 \tan \gamma}\left\{\frac{1}{x^4}\left(\frac{d^2 Z}{d x^2}+2 m \frac{d Z}{d x}+m^2 Z\right)-\frac{1}{x^6}\left(\frac{d Z}{d x}+m Z\right)\right\}\\
-\frac{4 n}{3 \tan \gamma} \frac{1}{x^5}\left(\frac{d Z}{d x}+m Z\right)
-2 n \frac{n^2-\cos 2 \gamma}{\sin \gamma \cos \gamma} \frac{1}{x^2} Z=0, \ldots . .(162)
\end{array}
\]
\todo[inline]{confirm that equation 163 is the third of triad of equations discussed
in this section.}
\[
\begin{array}{c}
\left(\frac{d X}{d x}+m X\right)+(2-\sigma) \frac{\cos \gamma}{x} X-2 n \frac{2-\sigma}{3} \frac{1}{x^2}
\left(\frac{d^2 Y}{d x^2}+2 m \frac{d Y}{d x}+m^2 Y\right) \\
+ 2n \frac{8-\sigma}{3} \frac{1}{x^3} \left(\frac{d Y}{d x}+m Y\right)+
\left( \frac{8 n^3}{3 \sin ^2 \gamma} -\frac{32}{3} n\right) \frac{Y}{x^4}+2 n Y \\
-\frac{\sin \gamma}{6}
\left\{\left(
\frac{d^4 Z}{d x^4}+4 m \frac{d^3 Z}{d x^3}+
6 m^2 \frac{d^2 Z}{d x^2}+4 m^3 \frac{d Z}{d x}+m^4 Z
\right)\right. \\
-\frac{2}{x}\left(\frac{d^3 Z}{d x^3}+3 m \frac{d^2 Z}{d x^2}+3 m^2 \frac{d Z}{d x}+m^3 Z\right)
-\frac{1}{x^2}\left( \frac{d^2 Z}{d x^2} +2 m \frac{d Z}{d x}+m^2 Z\right)+\\
\left.  -\frac{1}{x^2}\left( \frac{d^2 Z}{d x^2} +2 m \frac{d Z}{d x}+m^2 Z\right)+
 \frac{9}{x^3}\left(\frac{d Z}{d x}+m Z\right)\right\} \\
+\frac{2 \sigma n^2}{3 \sin \gamma}
\left\{\frac{1}{x^2}\left(\frac{d^2 Z}{d x^2}+2 m \frac{d Z}{d x}+m^2 Z\right)
-\frac{1}{x^3}\left(\frac{d Z}{d x}+m Z\right)\right\}+\\
+\frac{4 n^2}{3 \sin \gamma} \frac{1}{x^3}\left(\frac{d Z}{d x}+m Z\right) +
2 \frac{n^2-\cos^2 \gamma}{\sin \gamma} Z=0 .
\end{array}\]

\subsection{Solution of the equations.(c)}

Ve wish tc obtain solutions of these equations in series proceeding by
powers of \(x^{-1}\). write

\[E=\frac{1}{x},\]

issume for \(X, Y, Z\) expressions of the form

\[
\begin{aligned}
X&=z^a\left(a_0+a_1 s+a_2 z^2+\ldots\right)\\
Y&=z^b\left(b_0+b_1 z+b_2 z^2+\ldots\right)\\
Z&=z^e\left(c_0+c_1 s+c_2 z^2+\ldots\right)
\end{aligned}
\]
Then we have, for example,
\[
\frac{dZ}{dx} = -z^2 \frac{dZ}{dz}=-z^{c+1}
(c_o c + c_1(c+1)z + c_2(c+2)z^2 +\ldots),\nonumber \\
\]
\[
\frac{dZ^2}{dx^2} = -z^2 \frac{d}{dz}\frac{dZ}{dx}=
\]
\begin{equation}\label{eq164}
-z^{c+2}
(c_o c (c+1)+ c_1(c+1)(c+2)z + c_2(c+2)(c+3)z^2 + \ldots)
\end{equation}
and so on. We substitute the serial expressions for
\({\{\frac{d^k}{dx^k}\}}_{k=0}^4 Z\) in the left-hand
members of equations~\ref{eq161},\ref{eq162},\ref{eq163} and equate to
zero the coefficients of the various powers of \(z\),
beginning with the lowest.
We thus obtain sufficient equations to determine
\(m\), the exponents \(a\), \(b\), \(c\),
and the coefficients \({\{a_k\}}_{k\geq 0}\),
\({\{b_k\}}_{k\geq 0}\), \({\{c_k\}}_{k\geq 1}\),in terms of
\(c_o\), which remains arbitrary.
It will appear that there are four values of \(m\),
and so there are four solutions of this type.
If this process is carried out it will very
soon appear that
\(a=b=c=1/2\).  So it is simpler to assume the forms
\begin{equation}\label{eq165}
    (X,Y,Z)= z^{-1/2} \sum_{k \geq 0} (a_k, b_k, c_k ) z^k.
\end{equation}

We then have, for example,
\begin{align*}
\frac{dZ}{dx} &= 2^{-1} z^{\frac{1}{2} }
(c_o - c_1 z - 3 c_2 z^2 - 5 c_3 z^3 \ldots),\\
\frac{dZ^2}{dx^2} &= 2^{-2}z^{\frac{3}{2} }
(-c_o + 3 c_1 z + 5!! ~c_2 z^2 + \frac{7!!}{3} c_3 z^3 + \ldots),\\
\frac{dZ^3}{dx^3} &= 2^{-3} z^{\frac{5}{2} }
(3c_o - 5!! ~c_1 z - 7!! ~c_2 z^2 + \frac{9!!}{3} c_3 z^3 + \ldots),\\
\frac{dZ^4}{dx^4} &= 2^{-4}z^{\frac{7}{2} }
(- 5!! ~c_o + 7!! ~c_1 z - 9!! ~c_2 z^2 + \frac{11!!}{3} c_3 z^3 +\ldots),\\
\end{align*}
The terms of lowest degree in the three
equations are the terms containing \(z^{-1/2}\).
Equating their coefficients to zero,
we obtain the equations
\begin{align}
\frac{\sin \gamma }{2} m^2 a_o + \frac{1+ \sigma }{2}
\tan(\gamma) m~b_o+ \left( n^2 \frac{1+ \sigma }{2 \cos(\gamma) }
- \sigma \cos(\gamma) \right) m c_o & = 0 \nonumber \\
\frac{1-\sigma }{4} m^2 a_o + \frac{\sin^2 \gamma }{\cos(\gamma) } m^2 b_o+
\frac{1-\sigma }{4} n \tan(\gamma) m^2 c_o & = 0 \nonumber \\
\sigma \cos(\gamma) m~a_o + 2 m~b_o+
\left(
-\frac{\sin \gamma}{6} + 2 \frac{n^2 - \cos^2 \gamma }{\sin \gamma }
\right)c_o & = 0 \label{eq166}
\end{align}
yielding the results
\begin{equation}\label{eq167}
(a_o,b_o) = \left(\frac{2 \sigma }{m \tan(\gamma)}, -\frac{n}{\sin(\gamma)}\right) c_o,
\end{equation}
and the secular equation
\begin{equation}\label{eq168}
m^4 +12(1-\sigma^2) \cot^2(\gamma) =0.
\end{equation}
\todo[inline]{equation 169 is supposed to be equation 168} 

The four complex values of \(m\) given by (168)
are the same as occurred in the problem of
the conical shell strained by lateral forces,
which was discussed in Article 349. The
present solution applies to the case of
lateral forces, \(n=1\), as well as to conditions
answering to greater values of \(n\).

When \(m\) satisfies equation~(\ref{eq168}), and the left-hand
members of equations (166) are multiplied in order by
\setcounter{equation}{0}
\begin{equation}\label{coefficient_equation_168}
2 m \sigma \sin(\gamma),\quad
4( 2 + \sigma )n,\quad
-m^2 \sin(\gamma) \tan(\gamma),
\end{equation}
and added, the coefficients of \(a_o,b_o,c_o\) vanishes. This remark is
important because the same thing would happen if we formed the equations
like~(\ref{eq166}),
containing no coefficients \(a_k\), \(b_k\), \(c_k\) of suffix greater than
\(k\), by equating to zero the coefficients of \(z^{k-1/2}\),
in equations~\ref{eq161},\ref{eq162},\ref{eq163}, and
multiplied their left-hand members in order
by the same three expressions~(\ref{coefficient_equation_168})
then we should find that
\(a_k\), \(b_k\), \(c_k\)
would all disappear from the equation so formed. It is
therefore convenient to form a new equation from
equations~\ref{eq161},\ref{eq162},\ref{eq163} by multiplying
\todo[inline]{Page header 352. EDGE-EFFECT. UNSYMMETRICAL CONDITIONS 611}
\todo[inline]{The next equation is contains misprints}
the left-hand members of these equations in order by
\(2 \sigma m \sin \gamma, 4(2+\sigma) n\), and
\(-m^2 \sin \gamma \tan \gamma\) and adding. We thus obtain the equation

\[\begin{aligned}
& \sigma m \sin ^2 \gamma\left\{\left(\frac{d^2 X}{d x^2}+2 m \frac{d X}{d x}+m^2 X\right)-\frac{1}{x}\left(\frac{d X}{d x}+m X\right)-\frac{3}{x^2} X\right\}-2 m n^2 \sigma(1-\sigma) \frac{X}{x^2} \\
& -2 n^2(2+\sigma)\left\{\frac{1+\sigma}{x^2}\left(\frac{d X}{d x}+m X\right)+\frac{5-3 \sigma}{x^3} X\right\}-m^2 \sin ^2 \gamma\left\{\sigma\left(\frac{d X}{d x}+m X\right)+(2-\sigma) \frac{X}{x}\right\} \\
& +2 \sigma m \sin \gamma\left[n \frac{1+\sigma}{2} \tan \gamma\left(\frac{d Y}{d x}+m Y\right)-n(3-\sigma) \tan \gamma \frac{Y}{x}\right. \\
& 1-\sigma\left\{\begin{array}{c}
3 \tan \gamma \\
\left.\left.\frac{1}{x^4}\left(\frac{d Y}{d x}+m Y\right)-\frac{2}{x^6} Y\right\}\right]
\end{array}\right. \\
& \left.\left.+4(2+\sigma) n\left[\frac{1-\sigma \sin ^2 \gamma}{4}\right) \frac{\left(d^2 Y\right.}{\cos \gamma}+2 m \frac{d Y}{d x}+m^2 Y\right)+\frac{1}{x}\left(\frac{d Y}{d x}+m Y\right)-\frac{4}{x^2} Y\right\}-\frac{2 n^2}{\cos \gamma} \frac{Y}{x^2} \\
& \left.+(1-\sigma) \cos \gamma\left\{\frac{1}{x^4}\left(\frac{d^2 Y}{d x^2}+2 m \frac{d Y}{d x}+m^2 Y\right)-\frac{3}{x^5}\left(\frac{d Y}{d x}+m Y\right)+\frac{4}{x^6} Y\right\}-\frac{8 n^3 \cos \gamma}{3 \sin ^2 \gamma} \frac{Y}{x^3}\right] \\
& -m^2 \sin \gamma \tan \gamma\left[-2 n \frac{2-\sigma}{3} \frac{1}{x^2}\left(\frac{d^2 Y}{d x^2}+2 m \frac{d Y}{d x}+m^2 Y\right)+2 n \frac{8-\sigma}{3} \frac{1}{x^3}\left(\frac{d Y}{d x}+m Y\right)\right. \\
& +2 \sigma m \sin \gamma\left[\left(n^2 \frac{1+\sigma}{2 \cos \gamma}-\sigma \cos \gamma\right)\left(\frac{d Z}{d x}+m Z\right)-\left(n^2 \frac{3-\sigma}{\cos ^2 \gamma}-\frac{32 n}{3}\right) \frac{Y}{x^4}+2 n Y\right] \\
& +4(2+\sigma) n\left[\frac{1-\sigma}{4} n \tan \gamma\left\{\left(\frac{d^2 Z}{d x^2}+2 m \frac{d Z}{d x}+m^2 Z\right)+\frac{1}{x}\left(\frac{d Z}{d x}+m Z\right)-\frac{4}{x^2} Z\right\}\right. \\
& -2 n n^2-\cos ^2 \gamma Z \\
& \left.\sin \gamma \cos \gamma \frac{2 \sigma n}{x^2}-\frac{1}{3 \tan \gamma}\left\{\frac{d^2 Z}{x^4}+2 m \frac{d Z}{d x}+m^2 Z\right)-\frac{1}{x^6}\left(\frac{d Z}{d x}+m Z\right)\right\} \\
& -m^2 \sin \gamma \tan \gamma\left[-\frac{\sin \gamma}{6}\left\{\left(\frac{d^4 Z}{d x^4}+4 m \frac{d^3 Z}{d x^3}+6 m^2 \frac{d^2 Z}{d x}+4 m^3 \frac{d Z}{d x}+m^4 Z\right)\right.\right. \\
& \left.-\frac{2}{x}\left(\frac{d^3 Z}{d x^3}+3 m \frac{d^2 Z}{d x^2}+3 m^2 \frac{d Z}{d x}+m^3 Z\right)-\frac{1}{x^2}\left(\frac{d^2 Z}{d x^2}+2 m \frac{d Z}{d x}+m^2 Z\right)+\frac{9}{x^3}\left(\frac{d Z}{d x}+m Z\right)\right\} \\
& +\frac{2 \sigma n^2}{3 \sin \gamma}\left\{\frac{1}{x^3}\left(\frac{d^2 Z}{d x^2}+2 m \frac{d Z}{d x}+m^2 Z\right)-\frac{1}{x^3}\left(\frac{d Z}{d x}+m Z\right)\right\}+3 \sin ^2 \frac{1}{x^3}\left(\frac{d Z}{d x}+m Z\right) \\
& \left.+2 \frac{n^2-\cos 2}{d x} Z\right]=0
\end{aligned} \]

(169) In future we shall use equation (169) instead of equation (163). If from
equations (161), (162), (169) we select the terms containing
\(z^{\kappa-\frac{1}{2}}\), and equate them to zero, the first two will
contain \(a_\kappa, b_\kappa, c_\kappa\), but the third will not contain
any \(a\) 's, \(b\) 's or \(c\) 's with suffixes higher than
\(\kappa-1\). For example, taking \(\kappa=1\), equation (169) will
yield a new equation connecting \(a_0, b_0, c_0\), with which we shall
deal presently, and equations (161) and (162) will yield two equations
containing \(a_1, b_1, c_1\). Again, taking \(\kappa=2\), equation (169)
will yield a new equation containing \(a_1, b_1, c_1\), and equations
(161) and (162) will yield two equations containing \(a_2, b_2, c_2\).
Thus the equations, by which \(a_k, b_k, c_k\) are to be determined, are
obtained by equating to zero the coefficients of \(z^{x-i}\) in (161)
and (162), and the coefficient of \(z^{k+\frac{1}{2}}\) in (169).

612 CONICAL SHELL CH. XXIV A

In particular by equating to zero the coefficient of \(z^{\frac{1}{2}}\)
in equation (169) we obtain the equation

\[\begin{array}{c}
-\frac{1}{2} \sigma m^2 \sin ^2 \gamma \cdot a_0-m^2 \sin ^2 \gamma(2-\sigma) a_0+\frac{1}{2} \sigma(1+\sigma) m n \frac{\sin ^2 \gamma}{\cos \gamma} b_0+\\
-2 \sigma(3-\sigma) m n \frac{\sin ^2 \gamma}{\cos \gamma} b_0 \\
\quad+2(2+\sigma)(1-\sigma) m n \frac{\sin ^2 \gamma}{\cos \gamma} b_0+\frac{1}{2} \sigma m \sin \gamma\left(n^2 \frac{1+\sigma}{\cos \gamma}-2 \sigma \cos \gamma\right) c_0 \\
\quad-2 \sigma m \sin \gamma\left(n^2 \frac{3-\sigma}{\cos \gamma}-2 \cos \gamma\right) c_0+2 m n^2(2+\sigma)(1-\sigma) \tan \gamma \cdot c_0=0 .
\end{array}\]

On substituting for \(a_0\) and \(b_0\) from equations (167) this
equation becomes an identity. It may by observed that, if we had not
assumed the exponents \(a, b, c\) to be equal to \(-\frac{1}{2}\), but
had left them undetermined, it would still have been necessary that the
coefficient of \(z^{e+1}\) in an equation formed in the same way as
(169) should vanish identically, and this condition would have led to
the determination of \(a, b\), and \(c\).

We now equate to zero the coefficients of \(z^{\frac{1}{2}}\) in
equations (161) and (162), obtaining the equations

\[\begin{array}{c}
\frac{\sin \gamma}{2} m^2 a_1+\frac{1+\sigma}{2} n \tan \gamma\left(m b_1+\frac{1}{2} b_0\right)-n(3-\sigma) \tan \gamma \cdot b_0+\\
+\left(n^2 \frac{1+\sigma}{2 \cos \gamma}-\sigma \cos \gamma\right)\left(m c_1+\frac{1}{2} c_0\right)
-\left(n^2 \frac{3-\sigma}{\cos \gamma}-2 \cos \gamma\right) c_0=0
\end{array}\]

and

\[\frac{1-\sigma}{4} \frac{\sin ^2 \gamma}{\cos \gamma}\left(m^2 b_1+2 m b_0\right)+\frac{1-\sigma}{4} n \tan \gamma\left(m^2 c_1+2 m c_0\right)=0,\]

which, combined with (167), yield

\[a_1=\frac{2 \sigma}{m \tan \gamma} c_1-\frac{4-\sigma}{m^2 \tan \gamma} c_0, \quad b_1=-\frac{n}{\sin \gamma} c_1 .\]

We also equate to zero the coefficient of \(z^{\frac{3}{2}}\) in
equation (169), obtaining the equation

DMD: This looks wrong.  I do not want to touch it.

\[\begin{aligned}
& -\sigma m \sin ^2 \gamma\left(2 m a_1+\frac{15}{4} a_0\right) \cdot 2 n^2(2+\sigma)(1+\sigma) m a_0-2 n^2 \sigma(1-\sigma) m a_0 \\
& -m^2 \sin ^2 \gamma\left\{-\frac{1}{2} \sigma a_1+(2-\sigma) a_1\right\} \\
& \begin{array}{r}
-\frac{1}{2} \sigma(1+\sigma) m n \frac{\sin ^2 \gamma}{\cos \gamma} b_1-2 \sigma(3-\sigma) m n \frac{\sin ^2 \gamma}{\cos \gamma} b_1-(2+\sigma)(1-\sigma) n \frac{\sin ^2 \gamma}{\cos \gamma} \frac{15}{4} b_0 \\
-\frac{8 n^3(2+\sigma)}{\cos \gamma} b_0+2 n \frac{2-\sigma}{3} m^4 \sin ^2 \gamma b_0 \\
\cos \gamma
\end{array} \\
& -\sigma m \sin \gamma\left(n^2 \frac{1+\sigma}{2 \cos \gamma}-\sigma \cos \gamma\right) c_1-2 \sigma m \sin \gamma\left(n^2 \frac{3-\sigma}{\cos \gamma}-2 \cos \gamma\right) \\
& -(2+\sigma)(1-\sigma) n^2 \tan \gamma \frac{15}{4} c_0-8(2+\sigma) n^2 \frac{n^2-\cos { }^2 \gamma}{\sin \gamma \cos \gamma} c_0+\\
& -m^4 \frac{\sin ^3 \gamma}{3 \cos \gamma}\left(2 m c_1+\frac{11}{4} c_0\right) \\
& 2 \sigma n^2 m^4 \tan \gamma \cdot c_0=0,
\end{aligned}\]

and, on substituting from (167) and (170), this equation can be
expressed as a relation between \(c_1\) and \(c_0\), which yields, after
some reduction involving the use of (168),

\[l_1=-\left(\frac{4 n^2}{\sin ^3 y}+\frac{19}{8}\right) \frac{c_0}{m} .\]

From this value of \(c_1\) we find instead of (170)

\[a_1=-\left(\frac{8 n^2 \sigma}{\sin ^2 \gamma}+4+\frac{15}{4} \sigma\right) \frac{c_0}{m^2 \tan \gamma}, \quad b_1=\left(\frac{4 n^2}{\sin ^3 \gamma}+\frac{19 n}{8 \sin \gamma}\right) \frac{c_0}{m} .\]
